\chapter{Room 315 Rail Transport}
\label{ch:rail}
\chapteroverview{Explain the rail model from its configuration inputs and
coordinate transforms through routing, devices, shuttle lifecycle, motion
guards, and simulation-only presentation.}{Changing rail geometry or devices,
diagnosing route or pose disagreement, extending shuttle behavior, or verifying
which state is authoritative.}

\section{Chapter Source Map}

\scriptsize
\begin{longtable}{L{0.23\textwidth} L{0.69\textwidth}}
\caption{Authoritative sources for the Room 315 rail transport system}
\label{tab:rail-source-map}\\
\toprule
\tableheader Topic & Authoritative repository path(s) \\
\midrule
\endfirsthead
\toprule
\tableheader Topic & Authoritative repository path(s) \\
\midrule
\endhead
Rail launch and runtime integration & \file{mfja_robot_control_config/launch/room_315_dual_kinematic_shuttles.launch.py}; \file{mfja_robot_control_config/scripts/room_315_kinematic_shuttle_node.py}; \file{mfja_robot_control_config/scripts/room_315_kinematic_shuttle.py} \\ \rowseparator
Directed network and sampled geometry & \file{mfja_robot_control_config/config/room_315_kinematics/rail_network_right.yaml}; \file{mfja_robot_control_config/config/room_315_kinematics/rail_network_left.yaml}; \file{mfja_robot_control_config/config/room_315_kinematics/raw_segments/}; \file{mfja_robot_control_config/scripts/room_315_kinematic_shuttle.py} \\ \rowseparator
Calibration and public-name mapping & \file{mfja_robot_control_config/scripts/room_315_rail_defaults.py}; \file{mfja_robot_control_config/scripts/room_315_device_position_tool.py}; \file{mfja_robot_control_config/scripts/room_315_kinematic_shuttle_node.py} \\ \rowseparator
Slots, switches, stoppers, and sensors & \file{mfja_robot_control_config/config/room_315_kinematics/rail_devices_right.yaml}; \file{mfja_robot_control_config/config/room_315_kinematics/rail_devices_left.yaml}; \file{mfja_robot_control_config/scripts/room_315_rail_devices.py}; \file{mfja_robot_control_config/scripts/conveyor_loop_mode_controller.py} \\ \rowseparator
Identity and lifecycle & \file{mfja_robot_control_config/scripts/room_315_multi_shuttle.py}; \file{mfja_robot_control_config/scripts/room_315_kinematic_shuttle_node.py}; \file{mfja_rail_interfaces/srv/AddShuttle.srv}; \file{mfja_rail_interfaces/msg/ShuttleState.msg} \\ \rowseparator
Commands, states, and headway & \file{mfja_robot_control_config/scripts/room_315_kinematic_shuttle.py}; \file{mfja_robot_control_config/scripts/room_315_kinematic_shuttle_node.py}; \file{mfja_rail_interfaces/msg/ShuttleCommand.msg}; \file{mfja_3rd_floor_bringup/launch/room_315_floor_common.py} \\ \rowseparator
Markers and payload visuals & \file{mfja_robot_control_config/scripts/room_315_kinematic_shuttle_node.py}; \file{mfja_robot_control_config/scripts/room_315_rail_defaults.py}; \file{mfja_3rd_floor_description/models/room315_vla_overhead_devices/model.sdf}; \file{mfja_3rd_floor_description/models/room315_vla_payload_medium_box/model.sdf}; \file{mfja_3rd_floor_description/models/room315_vla_removable_obstacle_marker/model.sdf} \\
\bottomrule
\end{longtable}
\normalsize

\section{System Model and Authorities}

\subsection{Modelling decision}

Shuttle propulsion is kinematic. Each rail node loads directed topology and
calibrated CSV samples, advances a logical shuttle by arc length, transforms the
logical pose into the Gazebo frame, and calls the world \code{SetEntityPose}
service. No wheel torque or rail-contact dynamics propel a shuttle.

This design makes routing, devices, datasets, and effect checks deterministic,
but it has a crucial consequence: logical rail state may continue advancing if
a pose service is missing or times out. Monitor diagnostics and Gazebo together
with ROS state when visual/logical agreement matters.

\subsection{Authoritative input files}

Each rail side loads five source layers in the following order:

\begin{enumerate}
  \item \file{mfja_robot_control_config/config/room_315_kinematics/rail_network_<side>.yaml}:
        nodes, 14 forward-only segments, switches, fixed transitions, and the
        routing table;
  \item \file{mfja_robot_control_config/config/room_315_kinematics/raw_segments/*.csv}:
        \code{index,x,y,z} samples for each public segment filename;
  \item \file{mfja_robot_control_config/config/room_315_kinematics/rail_devices_<side>.yaml}:
        slots, binary sensors, stoppers, radii, and default states;
  \item \file{mfja_robot_control_config/scripts/room_315_rail_defaults.py}:
        public topics, identity/entity defaults, left-name mapping, Gazebo
        calibration, and marker styles; and
  \item identity and model conventions from
        \file{mfja_robot_control_config/scripts/room_315_multi_shuttle.py} and
        the identity-bearing shuttle models in
        \file{mfja_3rd_floor_description/models/}.
\end{enumerate}

The explicit \code{csv: raw_segments/<segment>.csv} field is authoritative.
Do not derive filenames from keys even though normalized live files currently
match. A narrow compatibility table exists only for immutable legacy snapshots.
The loaders are
\file{mfja_robot_control_config/scripts/room_315_kinematic_shuttle.py},
\file{mfja_robot_control_config/scripts/room_315_rail_devices.py}, and
\file{mfja_robot_control_config/scripts/room_315_rail_defaults.py}.

\section{Static Rail Definition}

\subsection{Geometry and path sampling}

\file{mfja_robot_control_config/scripts/room_315_kinematic_shuttle.py} provides
two backends:

\begin{itemize}
  \item \code{cubic_hermite} (default): continuous interpolation from sampled
        positions and tangents, reparameterized by numerically sampled arc
        length; and
  \item \code{polyline}: direct piecewise-linear interpolation for comparison
        and calibration debugging.
\end{itemize}

Each segment stores length $L$; devices use normalized
$s_{ratio}\in[0,1]$ and resolve $s=s_{ratio}L$. Motion advances
$\Delta s=v\Delta t$, consuming residual distance across valid successor
segments. Every segment follows CSV index 0 to the final row. Reverse traversal
is not supported.

\subsection{Calibration into Gazebo coordinates}

Logical CSV coordinates are first transformed by a side-specific affine map,
then scaled around a configured origin, rotated around a configured origin, and
offset. In compact form:

\[
\begin{bmatrix}x_b\\y_b\end{bmatrix}=
\begin{bmatrix}a&b\\c&d\end{bmatrix}
\begin{bmatrix}x\\y\end{bmatrix}+
\begin{bmatrix}t_x\\t_y\end{bmatrix},
\qquad
\mathbf{p}_{gz}=R(\theta)\,S(\mathbf{p}_b-\mathbf{o}_s)
+\mathbf{o}_r+\mathbf{t}.
\]

Exact coefficients live in
\file{mfja_robot_control_config/scripts/room_315_rail_defaults.py}. Runtime
scale, rotation, and offsets are node parameters. Use
\file{mfja_robot_control_config/scripts/room_315_device_position_tool.py} to
convert a Gazebo point into a candidate \code{segment+s_ratio}; never tune
topology to compensate for a calibration error.

\subsection{Directed topology and switch guards}

The right public cycle has two guarded routing stages:

\begin{longtable}{L{0.16\textwidth} L{0.47\textwidth} L{0.29\textwidth}}
\caption{Right-rail directed routing stages and matching convergence guards}
\label{tab:right-rail-routing-stages}\\
\toprule
\tableheader Stage & Diverging path & Required convergence \\
\midrule
\endfirsthead
\toprule
\tableheader Stage & Diverging path & Required convergence \\
\midrule
\endhead
A3/A4 half & A23 $\rightarrow$ A3E/A3I $\rightarrow$ A34E/A34I & Matching A4 exterior/interior guard $\rightarrow$ A4E/A4I $\rightarrow$ A14 \\ \rowseparator
A1/A2 half & A14 $\rightarrow$ A1E/A1I $\rightarrow$ A12E/A12I & Matching A2 exterior/interior guard $\rightarrow$ A2E/A2I $\rightarrow$ A23 \\
\bottomrule
\end{longtable}

A diverging switch selects an outgoing branch. A converging gate must match the
branch on which the shuttle arrives. An unknown state, missing successor, or
mismatched guard has no valid successor and sets \state{FALLING}; the runtime
does not teleport or silently select another path.

\subsubsection{Left internal versus public names}

The left rail reuses the raw geometry but rotates its public station/switch
vocabulary. Keep conversion direction explicit:

\begin{longtable}{L{0.26\textwidth} L{0.26\textwidth} L{0.38\textwidth}}
\caption{Mapping between raw/internal and public left-rail names}
\label{tab:left-rail-name-mapping}\\
\toprule
\tableheader Internal/raw group & Public left group & Meaning \\
\midrule
\endfirsthead
\toprule
\tableheader Internal/raw group & Public left group & Meaning \\
\midrule
\endhead
A3E/A3I, switch at raw A3 & A1E/A1I, public A1 & First public left diverger \\ \rowseparator
A4E/A4I, raw A4 & A2E/A2I, public A2 & First public left converger \\ \rowseparator
A1E/A1I, raw A1 & A3E/A3I, public A3 & Second public left diverger \\ \rowseparator
A2E/A2I, raw A2 & A4E/A4I, public A4 & Second public left converger \\ \rowseparator
A34E/A34I & A12E/A12I & Upper/lower exterior/interior long branch mapping \\ \rowseparator
A12E/A12I & A34E/A34I & Opposite long branch mapping \\ \rowseparator
A14 & A23 & Trunk mapping \\ \rowseparator
A23 & A14 & Opposite trunk mapping \\
\bottomrule
\end{longtable}

Use \code{internal_rail_segment_name_to_public} when publishing internal state
and \code{public_rail_segment_name_to_internal} when resolving public input.
The current map happens to be an involution in several pairs; that does not make
the two functions interchangeable. The mapping authorities are
\file{mfja_robot_control_config/config/room_315_kinematics/rail_network_left.yaml}
and \file{mfja_robot_control_config/scripts/room_315_rail_defaults.py}.

\section{Devices and Occupancy}

\subsection{Slots and identity-bearing DZI sensors}

\begin{longtable}{L{0.12\textwidth} L{0.23\textwidth} L{0.20\textwidth} L{0.23\textwidth} L{0.14\textwidth}}
\caption{Configured slot and identity-bearing DZI sensor locations}
\label{tab:dzi-slot-locations}\\
\toprule
\tableheader Side & Slot/sensor & Raw segment & $s_{ratio}$ & Radius \\
\midrule
\endfirsthead
\toprule
\tableheader Side & Slot/sensor & Raw segment & $s_{ratio}$ & Radius \\
\midrule
\endhead
Right & slot 1 / DZI1R & A12E & 0.411866742 & 0.09 m \\ \rowseparator
Right & slot 2 / DZI2R & A12E & 0.653073633 & 0.09 m \\ \rowseparator
Right & slot 3 / DZI3R & A34E & 0.447469343 & 0.09 m \\ \rowseparator
Right & slot 4 / DZI4R & A34E & 0.683726992 & 0.09 m \\ \rowseparator
Left & slot 1 / DZI1L & raw A34E & 0.428330934 & 0.09 m \\ \rowseparator
Left & slot 2 / DZI2L & raw A34E & 0.674370792 & 0.09 m \\ \rowseparator
Left & slot 3 / DZI3L & raw A12E & 0.427586575 & 0.09 m \\ \rowseparator
Left & slot 4 / DZI4L & raw A12E & 0.668424664 & 0.09 m \\
\bottomrule
\end{longtable}

The device YAML also defines DA approach/branch sensors around every switch.
Single-point devices store \code{segment} and \code{s_ratio}; multi-point
devices store a \code{points} list. Each stopper has a derived
\code{A<n>_STOPPER_SENSOR} located 0.10 m upstream along each stopper point.
Only segment, ratio/points, and radius determine occupancy. Descriptive
\code{switch} or \code{branch} labels do not move a detector.

\subsection{Switches, stoppers, and delayed state}

Typed switch and stopper commands schedule actual-state changes at current
simulation time plus configured delay (0.3 s switch, 0.1 s stopper by high-level
default). State topics report the delayed actual value. Routing uses actual
switch state, not the requested command.

Stopper state \state{0} means open/pass; \state{1} means closed/stop. A closed
stopper clamps an enabled shuttle at its trigger and places it in
\state{WAITING}; it resumes when the guard clears. Sensor output is binary
occupancy, including swept crossing between ticks. The additional
\code{segment}, \code{s}, and \code{s_ratio} fields are privileged simulator
metadata for debugging/evaluation, not extra physical measurements.

Accepted delayed switch actual state is mirrored as a comma-separated
\code{selector=EXTERIOR/INTERIOR} string on
\rosname{/mfja/conveyor/switch_cmd}. The visual switch controller updates blade
models and publishes a latched summary on
\rosname{/mfja/conveyor/switch_states}. While a typed target is outstanding,
matching feedback is only an acknowledgement and conflicting feedback is
ignored as stale. With \code{sync_from_visual_switch_states=true}, however, an
unmatched visual summary can schedule a delayed update of the logical switch
registry and therefore affect routing. Treat this subscription as configured
state synchronization, not presentation-only telemetry; Chapter
\ref{ch:interfaces} documents the cross-boundary authority in detail.

\section{Shuttle Runtime}

\subsection{Identity and lifecycle}

Canonical identities are \code{R1--R4} and \code{L1--L4}, with Gazebo entities
\code{room315_right_shuttle_1--4} and
\code{room315_left_shuttle_1--4}. Startup selection modes are:

\begin{itemize}
  \item \code{count_prefix}: count $n$ selects the first $n$ identities; and
  \item \code{explicit}: exact comma-separated identities select an arbitrary
        subset and must agree with the declared count.
\end{itemize}

Initial start slots/positions must match the selected count and be unique when
occupancy rejection is enabled. Zero is a valid count and produces a heartbeat
with an empty name so the presence provider can distinguish a known-empty rail
from a missing or stale rail.

The \code{AddShuttle} service resolves an empty name to the next unused entity
and an empty slot to the next available configured slot. A requested speed that
compares greater than zero is used; otherwise the configured default is
substituted, with no separate finiteness validation. Duplicate entities,
invalid slots, and---when enabled---occupied start slots are rejected. A
non-preloaded model is spawned asynchronously; a successful service reply means
the logical shuttle was accepted, not that Gazebo has already rendered it.
\state{REMOVE} removes logical management immediately and requests asynchronous
Gazebo deletion; a failed delete may leave a visual orphan.

\subsection{State machine and commands}

The command-driven transitions are separated from autonomous guard and target
transitions to make command verification unambiguous.

\begin{longtable}{L{0.16\textwidth} L{0.24\textwidth} L{0.18\textwidth} L{0.34\textwidth}}
\caption{Command-driven shuttle transitions and required verification}
\label{tab:shuttle-command-transitions}\\
\toprule
\tableheader Command & Applicable state & Result & Verification and semantic \\
\midrule
\endfirsthead
\toprule
\tableheader Command & Applicable state & Result & Verification and semantic \\
\midrule
\endhead
\state{ON} & \state{DISABLED}, \state{WAITING}, or \state{FALLING} & \state{MOVING} with positive retained or supplied speed & Require a newer state; the current code can clear \state{FALLING} without homing \\ \rowseparator
\state{OFF} & Any managed mode & \state{DISABLED} & Clear the exact target, retain configured speed, and require newer explicit feedback \\ \rowseparator
\state{RESET} & Any managed mode with a free start slot & \state{MOVING}, \state{WAITING}, or \state{DISABLED} & Verify the start pose and resulting mode; enabled state and retained speed are preserved \\ \rowseparator
\state{REMOVE} & Any managed mode & Removed from registry & Check both logical registry/state and Gazebo entity deletion \\
\bottomrule
\end{longtable}

\begin{longtable}{L{0.20\textwidth} L{0.28\textwidth} L{0.18\textwidth} L{0.26\textwidth}}
\caption{Autonomous shuttle transitions caused by guards and targets}
\label{tab:shuttle-automatic-transitions}\\
\toprule
\tableheader Current state & Condition & Result & Semantic \\
\midrule
\endfirsthead
\toprule
\tableheader Current state & Condition & Result & Semantic \\
\midrule
\endhead
\state{MOVING} & Stopper or headway guard & \state{WAITING} & Hold position until the active guard clears \\ \rowseparator
\state{WAITING} & Guard clear with positive speed & \state{MOVING} & Resume the retained command \\ \rowseparator
\state{MOVING} & Exact target reached & \state{DISABLED} & Clamp to target and populate \code{reached_target_slot} \\ \rowseparator
\state{MOVING} & No valid guarded successor & \state{FALLING} & Stop core stepping; do not choose another successor \\
\bottomrule
\end{longtable}

\code{ShuttleState.speed} is the retained commanded travel setting, not
instantaneous velocity. It may remain non-zero in \state{DISABLED} or
\state{WAITING}. \code{reached_target_slot} is populated only when the
low-level controller clamps to an explicit target and disables the shuttle.

\subsection{Headway control and failure behavior}

Before accepting a motion step, the node predicts center separation among
managed shuttles and uses bounded search to clamp progress before the configured
minimum (0.33 m by default). A held shuttle reports \state{WAITING} and resumes
when clear. This is a software center-distance rule, not certified block
protection or complete rigid-body collision avoidance.

\begin{safetybox}[Fail-closed rail conditions]
Unknown or mismatched route state enters \state{FALLING} and stops core
stepping, but the current positive-speed \state{ON} path can clear that mode
without a reset. Treat that behavior as a safety defect and recover only
through a verified supervisor procedure. An occupied start can reject
add/reset; stale typed feedback must not be replaced with a fixed delay.
Low-level typed topics can bypass planning, so manual publishers remain an
expert interface.
\end{safetybox}

\section{Simulation-only Markers and Payloads}

When enabled, the rail node spawns collision-free visual markers gradually via
the world create service. Sensors are blue/inactive and green/active; stoppers
are amber/open and red/closed; switches use green/interior and orange/exterior
in debug mode; a shuttle is red in \state{FALLING}. Dynamic remove/recreate
refresh is disabled by default to avoid races with pending entity creation.

Payload command/state topics use JSON because payload visuals are privileged
simulation/dataset controls, not the public physical rail contract. A payload
box follows its shuttle with configured offsets. Spawn/delete failure does not
change logical payload state, and no payload physics or grasp is simulated.
